% MT037 / explanation — canonical editable content.
均方误差（mean squared error，MSE）把每条预测误差平方后再取平均。对\(N\)条验证记录，\(\mathrm{MSE}=\frac1N\sum_i[f(x_i)-y_i]^2\)；同一批记录上，数值越小表示平均平方误差越小。


均值基线忽略所有输入，无论晴天雨天、早晨晚上，都预测训练集的平均出租车使用量。线性模型在验证集上与它的MSE差不多，说明当前模型还没有把输入中的信息转成可推广的预测收益。


\textbf{建议一：补充或改善有用特征。} 使用量可能与节假日、公共交通状态有关；还可能是“雨天并且下班高峰”才特别高。可以加入有依据的额外特征或交互项，例如雨天指标乘晚高峰指标，让模型区分只下雨、只晚高峰与两者同时发生。


\textbf{建议二：允许非线性关系。} 若直接把小时数字当一维线性输入，模型只会认为时间每增加一小时，需求固定增加或减少。实际需求可能早晚有两个高峰；可用适当的时间编码、非线性特征，或树等非线性模型表达这种形状。然后在相同验证划分上比较MSE。


还应检查特征与目标是否按小时正确对齐、正则是否过强及优化是否成功。题目只给验证结果，不能直接推断训练代码算错。


选读：含截距的无正则OLS在\textbf{同一个训练集}上最优时，不会比均值基线的平方误差更大，因为\(w=0,b=\bar y\)本来就是它可选的一组参数；找到最优解至少不差于这一组。验证集是新数据，没有这个保证，所以题目现象并不矛盾。


\textbf{出处：} 2021B*期中 Q11，原卷第4页（10分）


\textbf{答案依据：} 2021B*期中 Q11 随卷／配套参考答案；条件与解释经复核。
